\section{Consequences, Limitations, and Conclusion}\label{sec:consequences}

On the nonlinear fleet, the collinear theory and the simulation agree on the magnitude and on its stiffness and mass dependence, but not on the distribution over pairs. The post hoc magnitudes, $0.148$ by intervention and $0.161$ by design normalization, lie within about 10\% of the committed $0.165$, above the background-subtracted bound $0.093$ and far below the impulsive $0.44$.

By Corollary~\ref{cor:threshold}, a snap unloads a neighbor at $\Tn$ once $\bar T_{\mathrm p}$ reaches $\Tn/\rho$: at pretension, $2.3T_0$ for the impulsive model, $6.1T_0$ and $6.8T_0$ for the committed and interventional magnitudes, and $10.8T_0$ for the background-subtracted bound. These are extrapolated scale estimates, not design bounds: the magnitude sample stops at $4T_0$, and zero tension censors the taut response. Neighbors of scored snaps carry per-cell median tensions of $1.1$--$2.0T_0$, so $\rho\approx0.15$ places isolated unloading near $7$--$14T_0$, several times the median parent peak of $2.5T_0$; yet one parent snap in five exceeds $7T_0$. Of the 10\,735 recorded unloadings in Fig.~\ref{fig:threshold}, 69\% struck neighbors at or above pretension and 60\% followed peaks above $7T_0$; the counterfactual counts, reported as untrusted, sharpen this to 79\% and 82\% of the 4336 unloadings for which the snap was a necessary cause (78\% and 76\% after reweighting for the re-integration sampling). Snap-driven unloading on this testbed thus comes mainly from the upper tail of snap peaks acting on normally loaded neighbors, not from neighbors near slack. The finite-contact law should not be read as bounding how readily slack spreads: the pre-registered test that turned $R_{ij}$ and the neighbor margins into predicted slack onsets failed, under-predicting them in all 16 cells by factors of $0.39$--$0.80$ ($0.65$--$0.98$ against a post hoc target of neighbors taut at the snap). Cable severance, which redistributes steady tension rather than transmitting a snap, is treated in \cite{hajieghrary2026passive}.

The conclusions have three limits. The exact theory is collinear with lumped cables, whereas distributed lines carry the snap as a wave, and whether lever arms and yaw can raise a pair above the collinear worst case is open. The planar testbed has linear drag and no hull contact, and the interventional estimator (whose integrator failed its pre-registered check), the design normalization, and the timing comparison are post hoc. Two mass ratios, $0.24$ and $0.50$, test the rising branch of Fig.~\ref{fig:massratio}; the interior maximum near $1.4$ is numerical, not measured or proved. Measuring it, and the pairwise distribution of $\rho$ with its upper tail, on the planar fleet is future work.

\noindent\textit{Data, code, and pre-registration.} The declaration record (SHA-256 prefix \texttt{fc26c18c14b97b04}) fixed the plan, cells, tests, pass criteria, and forecast before the campaign, pins the digests of the plan, forecast, theory document, modules, and inputs it uses, and was never edited. Of eight dated addenda, one records a reducer-crash fix, four record post hoc sensitivities, corrections, or reporting constraints that change no declared test, threshold, population, or rule, and three declare the fixed-dissipation control and the step-size check before their runs and the second mass ratio after a four-seed pipeline check. The declarations, addenda, run records, analysis code, and figure scripts are available from the authors.
